% -----------------------------------------------------------------------------------------------------
% Hypotheses A1 and A4 (Appendix A1, theory/A1_lower_bound.tex). Restated here, in this order: def:span, thm:LZ,
% thm:HXZ, prop:constants, prop:minimizer, prop:anyseq, thm:instance, cor:count, prop:sagamean, cor:saga, lem:overlap,
% prop:A4.
% Numbers quoted after prop:A4 (computed from eq:sagawindow and eq:lkwindow with rho replaced by 1/105 and
% C0 by 1 + 5Q/8): window of Figure A4b, n = 16, m = 256, eps = 1e-8, 65 <= kappa <= kappa_max = 49428 (2.88 decades);
% upper/lower ratios 7.76--7.84 (SAGA) and 88.6--114.9 (L-Katyusha); end-to-end slopes allowed [0.691, 1.310] and
% [-0.175, 1.216]. L-Katyusha row of the table: Corollary cor:lk-complexity of Appendix A3.
% -----------------------------------------------------------------------------------------------------
\subsection{Lower bounds and the hard instances}\label{sec:th-A1}

\paragraph{Oracle classes, known bounds and hypothesis A1.}
Costs are counted in component gradients, a full gradient costing $n$; Table~\ref{tab:th-complexity} collects the
guarantees of the methods compared in this report and the main lower bounds. The lower bounds of the literature differ
in how an algorithm may choose the indices of its queries and combine the answers; the following classes make these
assumptions exact.
\restate{def:span}
Lower bounds for (S1) therefore apply to SAGA, up to its initial sweep, but not to L-SVRG or L-Katyusha, whose full
gradients are $n$ queries at one point; those for (S2) apply to GD, AGD, SAGA, L-SVRG and L-Katyusha. On the chain
constructions below, a span method reveals at most one new coordinate per query, so that span-based bounds require a
dimension of the order of the number of queries they certify.

Hypothesis A1 (Section~\ref{sec:intro}) asserts the bound $\Omega((n+\sqrt{n\kappa})\log\frac1\varepsilon)$ for class
(S1) and a weakened form for arbitrary index sequences. Its first part is the theorem of Lan and Zhou, quoted with the
displays of arXiv:1507.02000v3, whose sum convention gives the same function and bounds as the journal version
\cite{lan2018optimal}. On their instance \eqref{eq:LZinstance}, $m$ components act on disjoint blocks of dimension
$\tilde n$, each a Nesterov chain $f_i$, convex with an $L_i$-Lipschitz gradient, $L_i=\mu(\mathcal Q-1)$, plus its own
regularizer; $n=m\tilde n$ is the total dimension, and each block has condition number at most $\mathcal Q$.
\restate{thm:LZ}
Written as an average of $m$ components, $m$ times the summands of \eqref{eq:LZinstance}, the instance has
$R^2=m\mu\mathcal Q$, so that $n_{\rm book}=m$, $\kappa=m\mathcal Q$ and, as $\mathcal C=m(\mathcal Q-1)$,
$m+\sqrt{m\mathcal C}\asymp n_{\rm book}+\sqrt{n_{\rm book}\kappa}$: this is the bound
$\Omega((n+\sqrt{n\kappa})\log\frac1\varepsilon)$ for class (S1), the first part of A1, on instances with $\kappa\ge n$
and a dimension that grows with the number of queries.

The second part holds only in a restricted sense. Han, Xie and Zhang \cite{han2024lower} prove a lower bound for a
class $\mathcal A$ that, by their account, contains SAGA, L-SVRG and L-Katyusha: at each step, an algorithm of
$\mathcal A$ samples one index i.i.d.\ from a fixed distribution, may query the proximal incremental first-order oracle (PIFO, which returns the
value, the gradient and a proximal step of a component) of the sampled component at all previous iterates, computes
the full gradient at the current point with an independent probability at most $c_0/n$, for a constant $c_0$ (the full
gradient at the initial point is free), and takes its iterate in the span of the previous iterates and of the answers,
projected on the domain. In their theorem, $L$ is the smoothness constant of the components (our $R^2$), and $R$ is not
ours: $2R$ bounds the diameter of the domain.
\restate{thm:HXZ}
For the methods of class $\mathcal A$, and, up to logarithmic factors, for the oblivious methods of part (c) of
\resref{thm:HLOY}, the lower bound is
$\Omega\big(n+(\frac{n}{1+(\ln(n/\kappa))_+}+\sqrt{n\kappa})\ln\frac1\varepsilon\big)$: the $\sqrt{n\kappa}$ term is
untouched, and full gradients can divide $n\log\frac1\varepsilon$ by at most $1+(\ln(n/\kappa))_+$, a gain only when
$\kappa\ll n$, never below the bare $n$; Prox-SVRG with the parameters of \resref*{thm:HLOY}(b) attains this bound when
$\kappa=O(n)$. Without these restrictions, only an additive $n$ is proved for randomized algorithms, while
deterministic algorithms, cyclic incremental methods included, can need
$\Omega(n\sqrt\kappa\log\frac{\varepsilon_0}{\kappa\varepsilon})$ when $\kappa>73$, by \resref{thm:WS}, which excludes
for them any gain over $n\log\frac1\varepsilon$ at small $\varepsilon$; whether randomly reshuffled span methods, or
deterministic ones when $\kappa\le73$, can beat $n\log\frac1\varepsilon$ is not settled by these sources.

\paragraph{A hard instance with a closed-form minimizer.}
The instance of \resref{thm:LZ} gives each component its own regularizer, which forces $\kappa\ge n$. The instance (S)
of Appendix~\ref{sec:lower_bound} keeps one Nesterov chain per component but shares the regularizer among all components,
as in \cite{agarwal2015lower}, so that every $\kappa>1$ is allowed. With $n$ blocks $\theta^{(i)}\in\R^m$ and
$A_m(\zeta)$ the $m\times m$ tridiagonal matrix with $-1$ off the diagonal and $2$ on it, except
$(A_m(\zeta))_{mm}=\zeta\in[1,3]$ (see \eqref{eq:chainmatrix}),
\begin{equation*}\tag{\ref*{eq:instance}}
\begin{aligned}
&c=\frac{\mu(\kappa-1)}{4},\qquad \phi(x)=\frac12\inner{A_m(\zeta)x}{x}-x_1\ \ (x\in\R^m),\\
&f_i(\theta)=c\,\phi(\theta^{(i)})+\frac\mu2\norm{\theta}^2,\quad i=1,\dots,n .
\end{aligned}
\end{equation*}
The instance (H) of Section~\ref{sec:setting}, used in Figures~\ref{fig:A3}, \ref{fig:A2} and~\ref{fig:A4}, has instead
a regularizer local to each block and $\kappa>n$. The constants and the minimizer of (S) are explicit.
\restate{prop:constants}
On (S), full-gradient methods therefore depend on $L_F\le\mu Q$, about $R^2/n$ when $\kappa\gg n$, rather than on $R^2$:
gradient descent and AGD do not pay there the factors $n$ and $\sqrt n$ of Table~\ref{tab:th-complexity}, which require
aligned components, as in the aligned instance (A) of Figure~\ref{fig:A4b}.
\restate{prop:minimizer}
Part (a) gives $\etastar$ in closed form, so that every error measured on (S) is exact; part (b), with $Q=1+4c/\mu_b$,
gives the minimizer of (H) (Section~\ref{sec:setting}).

Started from $\theta_0=0$, a span method reveals at most one new coordinate of block $j$ per query to component $j$
by \resref{lem:support}, and the coordinates not yet revealed carry an explicit share of $\norm{\etastar}^2$
by \resref{prop:tail}. Averaging over the blocks by Jensen's inequality gives a bound for every index sequence, the
$\sqrt{n\kappa}$ term of A1:
\restate{prop:anyseq}
For the methods with coin flips, such as L-SVRG and L-Katyusha, the bound holds for the expected number of gradient
evaluations (Appendix~\ref{sec:lower_bound}, after the proof). AGD on $F$, which needs $O(n\sqrt Q\log\frac1\varepsilon)$
gradients since $L_F\le\mu Q$, attains its order when $\kappa\gtrsim n$. The sequence-free argument cannot produce the
$n\log\frac1\varepsilon$ term, because a full gradient advances all $n$ chains at once. With i.i.d.\ indices, the number
of queries to each block is binomial, and the term appears:
\restate{thm:instance}
\restate{cor:count}
This proves the first part of A1 on (S), for every $\kappa>1$, under an explicit dimension condition: by \eqref{ineq:dim}
each block must be longer than the number of touches that an accelerated method needs, of order
$\frac14\sqrt{\kappa/n}\log\frac1\varepsilon$ for $\kappa\gg n$, so that the total dimension $nm$ is of the order of the
number of queries certified. L-Katyusha is not in (S1), because of its full gradients: its lower bound on (S) is
\resref{prop:anyseq}, without the $n\log\frac1\varepsilon$ term.

\paragraph{SAGA and L-Katyusha on the instance (hypothesis A4).}
Hypothesis A4 predicts that the gradient counts of SAGA and L-Katyusha grow like $\kappa$ and $\kappa^{1/2}$ when
$\kappa\gg n$. The chain bounds above are at most of order $(n+\sqrt{n\kappa})\log\frac1\varepsilon$; for SAGA,
the exponent $1$ comes instead from its mean iterate, which moves like gradient descent whatever the memory:
\restate{prop:sagamean}
For $n$ identical components this is the argument of \resref{prop:B3-lower}; here it holds on every quadratic instance.
Combined with the proof of Proposition~\ref{prop:59}, run with $\mu_F=\lmin(\nabla^2F)$ in place of $\mu$, it gives a
two-sided bound:
\restate{cor:saga}
On (S), the eigenvectors of the chain are sampled sine waves, and those of curvature at most $2\mu$ carry a fixed share of
the initial error as long as the chain is long compared with the decay length $\approx\sqrt Q/2$ of $\etastar$:
\restate{lem:overlap}
These results and the bound of Theorem~11 of \cite{kovalev2020dont} for L-Katyusha (Section~\ref{sec:th-A3} recalls
its iteration, its mirror sequence and its reference point, written $z_t$ and $w_t$ there and $z^k$ and $w^k$ below, and
its parameters $\theta_1$ and $\sigma$), converted into the distance of its mirror sequence in
\eqref{eq:LKconst}--\eqref{eq:LKupper} with an explicit constant $C_0\ge1$ that depends only on $n$, $\kappa$ and
$L_F/\mu$ and enters the counts only through $\log\frac{C_0}\varepsilon$, give the prediction tested in
Section~\ref{sec:exp-A4}:
\restate{prop:A4}
Part (ii) certifies the exponents $1$ and $\frac12$ as orders of growth, not as slopes. For the instance of
Figure~\ref{fig:A4b} ($n=16$, $m=256$, $\varepsilon=10^{-8}$) the window is $65\le\kappa\le\kappa_{\max}\approx4.9\cdot10^4$,
i.e.\ $2.9$ decades; with $\rho$ and $C_0$ replaced by their bounds $\frac1{105}$ and $1+\frac58Q$, the two sides of
\eqref{eq:sagawindow} differ there by a factor of about $7.8$, and those of \eqref{eq:lkwindow} by $88$ to $115$, so that
they constrain the end-to-end slope only to $[0.69,1.31]$ and $[-0.17,1.22]$, as \resref{rem:plot} details; the slopes
measured in Section~\ref{sec:exp-A4} agree with these orders but are not implied by them. Figure~\ref{fig:A4b} counts
L-Katyusha as in Table~\ref{tab:th-complexity} ($3$ gradients per iteration in expectation) and measures the distance
of its gradient sequence $y_t$: this multiplies the upper bound of \eqref{eq:lkwindow} by at most $\frac32$ and
replaces $C_0$ by a constant $C_y=O(\kappa)$ (Appendix~\ref{sec:lower_bound}, after \eqref{eq:LKupper}), without
changing the exponents.
Part (iii) proves that the counts saturate as $\kappa\to\infty$ only for SAGA and for L-Katyusha tuned with
$\lmin(\nabla^2F)$, the tuning of the experiments, and proves nothing of the kind with the nominal $\mu$; the slopes
measured beyond $\kappa_{\max}$ are not predicted by part (ii). Logistic regression is not a hard instance: there is no
chain, and its effective constants can be much better than $R^2$ and $\lambda$ (Section~\ref{sec:th-B2}).

\paragraph{Known and new.}
The lower bounds and the chain constructions \cite{nesterov2018lectures,agarwal2015lower,lan2018optimal} are known, and so
is the geometric minimizer given by $\zeta_Q$ \cite{lan2018optimal}. To our knowledge, the following are not in the
sources we reviewed in this form: the restatement of the bounds with their exact oracle classes and the correction of
hypothesis A1; self-contained proofs on a shared-regularizer instance valid for every $\kappa>1$, including the bound for
every index sequence (\resref*{prop:anyseq}); the explicit dimension condition and window (\resref*{cor:count},
\eqref{eq:window}); and the SAGA sandwich (\resref*{cor:saga}, with \resref*{lem:overlap}), which makes the exponent $1$
a theorem on the instance (\resref*{prop:A4}).

% Table 1 is anchored at the end of the subsection on purpose: in the first half of the subsection most pages begin
% inside a restated statement, so that a top float anchored there (at its first mention, or before the paragraph on
% the hard instance) interrupts a statement; anchored here, it lands among plain paragraphs (end of 4.4 or start of
% 4.5). Check the page of the table after any change of the layout.
\begin{table}[t]
\centering
\footnotesize
\setlength{\tabcolsep}{4pt}
\renewcommand{\arraystretch}{1.2}
\begin{tabular}{@{}p{3.0cm}p{6.6cm}p{6.3cm}@{}}
\toprule
Method or class & Component gradients to reach accuracy $\varepsilon$ & Validity and source \tabularnewline
\midrule
GD, step $1/R^2$ & $n\kappa\log\frac1\varepsilon$ \ ($n\frac{L_F}\mu\log\frac1\varepsilon$ with step $1/L_F$) & \raggedright
all $n,\kappa$; \cite[Section~5.4.4]{bach2024learning} \tabularnewline
AGD & $n\sqrt\kappa\log\frac1\varepsilon$ \ ($n\sqrt{L_F/\mu}\,\log\frac1\varepsilon$ with $L_F$) & \raggedright
all $n,\kappa$; \cite[Section~5.2.5]{bach2024learning} \tabularnewline
SAGA, $\gamma=\frac1{4R^2}$ & $n+\max\{3n,\frac{16\kappa}3\}\log\frac{1+n/4}{\varepsilon}$ & \raggedright all $n,\kappa$;
Proposition~\ref{prop:59}; \resref*{cor:saga} \tabularnewline
\raggedright L-SVRG, $\eta=\frac1{6R^2}$, $p=\frac1n$ & $n+3\max\{6\kappa,2n\}\log\frac{\Phi_0}{\varepsilon}$ &
\raggedright all $n,\kappa$; \cite[Theorem~5]{kovalev2020dont} \tabularnewline
L-Katyusha & $n+(12n+15\sqrt{n\kappa})\log\frac{3R^2(n+3)\norm{\theta_0-\etastar}^2}{8\varepsilon}$ & \raggedright
all $n,\kappa$ and $\varepsilon\le\frac{3R^2(n+3)}{16}\norm{\theta_0-\etastar}^2$; \resref*{cor:lk-complexity}
\tabularnewline
\midrule
class (S1) & $\Omega\big((n+\sqrt{n\kappa})\log\frac1\varepsilon\big)$ & \raggedright dimension $\gtrsim$ number of
queries; \resref*{thm:LZ}; \resref*{cor:count} \tabularnewline
\raggedright class (S2), instance (S) & $\frac{\sqrt{n(\kappa-1)}}8\log\frac1{2\varepsilon}$ & \raggedright
$\kappa-1\ge2n$, $q^{2m}\le\varepsilon<\frac12$; \resref*{prop:anyseq} \tabularnewline
\raggedright any randomized algorithm & $\Omega\big(n+\sqrt{n\kappa}\log(\frac{\varepsilon_0}{\varepsilon}\sqrt{n/\kappa})\big)$
& \raggedright $\kappa>161n$, no span assumption; \resref*{thm:WS} \tabularnewline
\bottomrule
\end{tabular}
\caption[Guarantees and lower bounds in component gradients]{Guarantees for every instance of the class of
Proposition~\ref{prop:59} (top) and lower bounds (bottom), in component gradients. The accuracy criteria differ only
inside the logarithm, by factors polynomial in $n$ and $\kappa$ (for instance $\frac{3\kappa(n+3)}4$ when the L-Katyusha
criterion is converted into the relative squared distance by strong convexity): relative squared distance (SAGA,
instance (S)), the Lyapunov function $\Phi$ of \cite[Theorem~5]{kovalev2020dont} (L-SVRG), $\E[F(y_t)]-F(\etastar)$
(L-Katyusha, $y_t$ its gradient sequence, Section~\ref{sec:th-A3}), initial gap $\varepsilon_0$ (last row). Counting convention, that of the experiments
(Section~\ref{sec:setting}): the SAGA, L-SVRG and L-Katyusha counts include the initial full gradient, and L-SVRG and
L-Katyusha are charged $2$ gradients per iteration ($f'_{i(t)}$ at the query point and at the reference point) and $n$
per refresh of the reference point, i.e.\ $3$ per iteration in expectation. With stored reference gradients, the
convention of $T_{\rm LK}$ in Appendix~\ref{sec:lower_bound}, an iteration costs $2$ in expectation, which gives
$n+(8n+10\sqrt{n\kappa})\log(\cdot)$ for L-Katyusha and $n+2\max\{6\kappa,2n\}\log(\cdot)$ for L-SVRG. The SAGA and
L-SVRG counts treat the number of iterations as a real number (rounding it up adds at most one iteration); the
L-Katyusha count is that of an integer number of iterations, which its condition on $\varepsilon$ ensures. For
$\kappa\gg n$ the leading terms are $n\kappa$, $n\sqrt\kappa$, $\kappa$, $\kappa$ and $\sqrt{n\kappa}$, the last one
matching the lower bounds.}
\label{tab:th-complexity}
\end{table}
